\documentclass{article}
\usepackage{graphicx} % Required for inserting images
\usepackage[T1]{fontenc}
\usepackage{amsmath}
\usepackage[margin=1in]{geometry}
\usepackage{booktabs}
\usepackage{tabularx}
\usepackage[hidelinks]{hyperref}

\title{Lesson 7: Forecasting Volatility}
\author{Analyst onboarding \textperiodcentered{} Lesson 7 of 10 \textperiodcentered{} L/S Gamma Desk, QUANTT}
\date{2026-10-02}

\begin{document}

\maketitle

\textbf{Goal:} understand how we forecast SPY's realized volatility over the life of a trade, and how to judge whether a forecast is any good.

\textbf{You need:} Lessons 1 and 6 (volatility patterns, RV, the VRP).

\section{Why we forecast}

IV is the market's forecast of future volatility. To trade against it, we need \textbf{our own forecast} of realized volatility over exactly the life of the option. If ours is better, the VRP (IV \ensuremath{-} our forecast) tells us which side to take.

Volatility \emph{can} be forecast because of the patterns in Lesson 1: \textbf{clustering} (today's volatility predicts tomorrow's), \textbf{mean reversion} (extremes fade toward normal), and the \textbf{leverage effect} (selloffs raise volatility).

\section{Model 1: the naive forecast}

``Next month will look like last month'': use the last 21 days' RV. It's simple, and surprisingly hard to beat by much. Treat it as the \textbf{baseline} every model must beat.

\section{Model 2: GARCH(1,1), our live model}

\textbf{GARCH} (Generalized AutoRegressive Conditional Heteroskedasticity) updates tomorrow's variance from three pieces:

\[\sigma^2_{t+1} = \omega + \alpha\, r_t^2 + \beta\, \sigma^2_t\]

\begin{itemize}

\item $\omega$: a constant that pulls volatility toward its \textbf{long-run level}.

\item $\alpha\, r_t^2$: \textbf{today's shock}. A big move today raises tomorrow's volatility.

\item $\beta\, \sigma^2_t$: \textbf{memory}. Today's volatility carries into tomorrow.

\end{itemize}

Typical values for SPY are $\alpha \approx 0.1$ and $\beta \approx 0.85$. Because $\alpha + \beta < 1$, a forecast drifts back toward the long-run variance $\omega / (1 - \alpha - \beta)$ the further ahead you look. That is mean reversion, built in.

The parameters are \textbf{fitted} to past returns by maximum likelihood: the computer picks the $\omega, \alpha, \beta$ that make the observed returns most probable.

\section{Model 3: HAR, volatility at three speeds}

The \textbf{HAR} model (Heterogeneous AutoRegressive) says different traders react on different time scales: day traders to yesterday, swing traders to last week, investors to last month. So tomorrow's RV is a weighted mix of three averages:

\[RV_{t+1} = c + \beta_d\, RV_{day} + \beta_w\, RV_{week} + \beta_m\, RV_{month}\]

The weights are fitted by ordinary least squares regression. HAR is simple and fast, and it usually beats GARCH when good RV data is available.

\section{Forecasting over the right horizon}

We don't need tomorrow's volatility. We need the \textbf{average volatility from today until the option expires}:

\begin{enumerate}

\item Count the trading days to the target expiry. 29 calendar days is about 20 trading days ($\approx 29 \times 252/365$).

\item Forecast the variance for each of those days. Each step leans further toward the long-run level.

\item Average the daily variances, annualize (\ensuremath{\times} 252) and take the square root.

\end{enumerate}

That annualized number is \texttt{rv\_\allowbreak{}forecast}, directly comparable with the option's IV.

\section{How to judge a forecast}

\begin{enumerate}

\item \textbf{Only out-of-sample.} Fit on the past, then forecast days the model never saw (Lesson 9). A model scored on its own training data always looks great.

\item \textbf{Use a fair loss.} Mean squared error is dominated by a few crash days. \textbf{QLIKE}, $\text{QLIKE} = \frac{RV}{\hat{\sigma}^2} - \ln\frac{RV}{\hat{\sigma}^2} - 1$ (with $RV$ in variance units), penalizes under-forecasting more and is the standard loss for volatility.

\item \textbf{Beat IV as a forecast.} IV itself is a forecast, though biased high by the VRP. A model worth trading should add information beyond IV.

\item \textbf{What finally matters is P\&L.} A forecast with slightly better accuracy but the same trading signals adds nothing. Judge by the trades it produces.

\end{enumerate}

\section{Beyond GARCH and HAR}

Our research (see \texttt{research/\allowbreak{}strategy\_\allowbreak{}research/\allowbreak{}}) has looked at:

\begin{itemize}

\item \textbf{Asymmetric models} (GJR-GARCH, leverage HAR) that let down moves raise volatility more than up moves.

\item \textbf{HAR with extra inputs} such as IV, the VIX term structure and recent momentum.

\item \textbf{Machine learning:} random forests, neural networks, and pretrained ``foundation'' time-series models such as TimesFM. They help modestly at best, and need careful out-of-sample testing.

\item \textbf{Combining forecasts:} averaging several models often beats choosing one. Lesson 10 covers using RL to choose the weights.

\end{itemize}

\section{In our code}

\begin{itemize}

\item \texttt{forecasting/\allowbreak{}GARCH.\allowbreak{}py} and \texttt{forecasting/\allowbreak{}HAR.\allowbreak{}py} share one interface: \texttt{fit(returns)}, \texttt{predict(horizon)} (returns annualized volatility) and \texttt{params()}.

\item \texttt{forecasting/\allowbreak{}\_\allowbreak{}\_\allowbreak{}init\_\allowbreak{}\_\allowbreak{}.\allowbreak{}py} has the registry: \texttt{FORECASTERS = \{``garch'': GARCH,\allowbreak{} ``har'': HAR\}}.

\item Switch models by editing one line in \texttt{configs/\allowbreak{}paper.\allowbreak{}toml}: \texttt{forecaster = ``har''}.

\item To add a model, write a class with the same three methods and register it.

\end{itemize}

\section{Words to know}

\begin{itemize}

\item \textbf{GARCH:} a volatility model driven by today's shock plus memory.

\item \textbf{HAR:} RV forecast from daily, weekly and monthly averages.

\item \textbf{Horizon:} how far ahead you forecast (for us, to expiry).

\item \textbf{Out-of-sample:} data the model was not fitted on.

\item \textbf{QLIKE:} the standard loss function for volatility forecasts.

\end{itemize}

\section{Check yourself}

\begin{enumerate}

\item After a \ensuremath{-}4\% day, what does GARCH do to tomorrow's forecast, and which term is responsible?

\item Why do we average the forecast over the days to expiry rather than use tomorrow's forecast?

\item Model A has lower forecast error than Model B but produces the same trades. Which should we use?

\end{enumerate}

\emph{Answers: (1) Raises it, through the $\alpha r_t^2$ shock term. (2) The option's payoff depends on volatility over its whole life. (3) Either: the trades are identical, so the extra accuracy is worth nothing. Keep the simpler one.}

\end{document}
